% TMS-001100.tex
% VSEPR-SIM v5.0.0-beta.6 — Formation and Premacro Analysis Foundation
% Release Report: Trajectory-Truth Analysis, Crystal Imperfections, Diffusion, and Macro Inference
% Branch: 4.0-legacy-beta

\documentclass[11pt,a4paper]{article}

% ---------------------------------------------------------------------------
% Packages
% ---------------------------------------------------------------------------
\usepackage[T1]{fontenc}
\usepackage[utf8]{inputenc}
\usepackage[margin=2.5cm]{geometry}
\usepackage{booktabs}
\usepackage{longtable}
\usepackage{array}
\usepackage{tabularx}
\usepackage{xcolor}
\usepackage{listings}
\usepackage{hyperref}
\usepackage{microtype}
\usepackage{parskip}
\usepackage{titlesec}
\usepackage{fancyhdr}
\usepackage{amsmath}
\usepackage{graphicx}
\usepackage{enumitem}
\usepackage{mdframed}

% ---------------------------------------------------------------------------
% Colours
% ---------------------------------------------------------------------------
\definecolor{vseprdark}{RGB}{30,40,60}
\definecolor{vsepraccentblue}{RGB}{50,100,180}
\definecolor{vsepraccentgold}{RGB}{180,140,40}
\definecolor{passgreen}{RGB}{30,130,60}
\definecolor{warnred}{RGB}{180,40,40}
\definecolor{codebg}{RGB}{245,245,248}
\definecolor{codefg}{RGB}{40,40,60}

% ---------------------------------------------------------------------------
% Hyperref
% ---------------------------------------------------------------------------
\hypersetup{
	colorlinks=true,
	linkcolor=vsepraccentblue,
	urlcolor=vsepraccentblue,
	citecolor=vsepraccentblue,
	pdftitle={TMS-001100 VSEPR-SIM v5.0.0-beta.6},
	pdfauthor={VSEPR-SIM Development},
	pdfsubject={Formation and Premacro Analysis Foundation}
}

% ---------------------------------------------------------------------------
% Headers / Footers
% ---------------------------------------------------------------------------
\pagestyle{fancy}
\fancyhf{}
\fancyhead[L]{\small\textcolor{vseprdark}{\textbf{VSEPR-SIM}} \textcolor{gray}{v5.0.0-beta.6}}
\fancyhead[R]{\small\textcolor{gray}{TMS-001100}}
\fancyfoot[C]{\small\textcolor{gray}{\thepage}}
\renewcommand{\headrulewidth}{0.4pt}

% ---------------------------------------------------------------------------
% Section formatting
% ---------------------------------------------------------------------------
\titleformat{\section}{\large\bfseries\color{vseprdark}}{}{0em}{\thesection\quad}[\titlerule]
\titleformat{\subsection}{\normalsize\bfseries\color{vsepraccentblue}}{}{0em}{\thesubsection\quad}
\titleformat{\subsubsection}{\normalsize\bfseries\color{vseprdark}}{}{0em}{\thesubsubsection\quad}
\titlespacing{\section}{0pt}{1.4em}{0.6em}
\titlespacing{\subsection}{0pt}{1em}{0.4em}

% ---------------------------------------------------------------------------
% Code listings
% ---------------------------------------------------------------------------
\lstset{
	backgroundcolor=\color{codebg},
	basicstyle=\ttfamily\small\color{codefg},
	frame=single,
	framerule=0.4pt,
	rulecolor=\color{gray!40},
	breaklines=true,
	keepspaces=true,
	showstringspaces=false,
	tabsize=4
}

% ---------------------------------------------------------------------------
% Status badge commands
% ---------------------------------------------------------------------------
\newcommand{\statpass}{\textbf{\textcolor{passgreen}{PASS}}}
\newcommand{\statwarn}{\textbf{\textcolor{warnred}{WARN}}}

% ---------------------------------------------------------------------------
% Framed note environment
% ---------------------------------------------------------------------------
\newmdenv[
	linecolor=vsepraccentgold,
	linewidth=1pt,
	backgroundcolor=vsepraccentgold!6,
	innerleftmargin=10pt,
	innerrightmargin=10pt,
	innertopmargin=6pt,
	innerbottommargin=6pt,
	skipabove=8pt,
	skipbelow=8pt
]{docnote}

% ---------------------------------------------------------------------------
% Title
% ---------------------------------------------------------------------------
\begin{document}

\begin{titlepage}
\vspace*{3cm}
{\color{vseprdark}\rule{\linewidth}{1.5pt}}\\[0.6em]
{\Huge\bfseries\color{vseprdark} TMS-001100}\\[0.3em]
{\large\color{vsepraccentblue} VSEPR-SIM v5.0.0-beta.6 --- Formation and Premacro Analysis Foundation}\\[0.2em]
{\normalsize\color{gray} Release Report: Trajectory-Truth Analysis, Crystal Imperfections, Diffusion, and Macro Inference}\\[0.4em]
{\color{vseprdark}\rule{\linewidth}{0.6pt}}

\vspace{2em}

\begin{tabular}{ll}
\textbf{Report ID}    & TMS-001100 \\[0.3em]
\textbf{Version}      & v5.0.0-beta.6 \\[0.3em]
\textbf{Branch}       & \texttt{4.0-legacy-beta} \\[0.3em]
\textbf{Commit label} & Formation and Premacro Analysis Foundation \\[0.3em]
\textbf{Date}         & 2025 \\[0.3em]
\textbf{Status}       & Instrumentation release complete; beta-7 target: unification \\
\end{tabular}

\vfill
{\color{gray}\small Generated from live test run. All numerical values confirmed reproducible across multiple executions.}
\end{titlepage}

% ---------------------------------------------------------------------------
\tableofcontents
\newpage

% ===========================================================================
\section{Abstract}
% ===========================================================================

VSEPR-SIM version 5 is the penultimate version iteration, finally touching on true macro-scale
properties. Beta-6 fills in several mechanisms for lower-level verification before doing
so --- implications initially skipped for simplification in previous versions but impossible to
ignore when moving onto very complex systems in coming developments (v5-beta.8).

\begin{docnote}
Beta-6 is the \textbf{instrumentation release}. It does not claim to be macro. It claims to
have built the instruments required before macro can be trusted.
\end{docnote}

The result is a working trajectory-to-inference pipeline that can measure structural change,
surface interaction, diffusion behavior, and bead/powder packing from trajectory data ---
\emph{without} storing derived material properties in simulation state. This distinction is the
entire project philosophy in practice.

% ===========================================================================
\section{Version Identity}
% ===========================================================================

\begin{table}[h!]
\centering
\begin{tabularx}{\linewidth}{llX}
\toprule
\textbf{Version} & \textbf{Role} & \textbf{Description} \\
\midrule
v4.1.0 & Stable foundation & Frozen known-good baseline \\
\textbf{v5.0.0-beta.6} & \textbf{Instrumentation} & Formation + premacro analysis foundations complete \\
v5.0.0-beta.7 & Unification & Formation outputs feed analyzers and reports \\
v5.0.0-beta.8 & Coupling & Analyzers feed macro-directed candidate workflows \\
\bottomrule
\end{tabularx}
\caption{Version roadmap context.}
\end{table}

Beta-6 was developed alongside v4.1.x but with entirely different goals. v4.1.x is a
stability-preserving branch. v5-beta is an architecture-expanding branch. They do not conflict.

% ===========================================================================
\section{Core Design Rule --- Preserved Across All New Systems}
% ===========================================================================

\begin{docnote}
\textbf{No inferred macro property is written back into State.}
\end{docnote}

The system respects four clean layers:

\begin{table}[h!]
\centering
\begin{tabularx}{\linewidth}{lX}
\toprule
\textbf{Layer} & \textbf{Role} \\
\midrule
\textbf{State}    & Position, velocity, identity, dynamic truth \\
\textbf{xyzFull}  & Trajectory evidence \\
\textbf{Analysis} & Measures behavior from trajectory \\
\textbf{Inference}& Estimates macro behavior from measured behavior \\
\bottomrule
\end{tabularx}
\caption{Strict four-layer architecture.}
\end{table}

Quantities explicitly \textbf{not} stored in State include: diffusion coefficient, effective
diffusivity by axis, transport class, porosity, bulk density, packing class, compressibility
proxy, permeability proxy, sintering readiness, and formation energy proxy.

These are analysis-only outputs. There is no field in \texttt{atomistic::State} for any of
them. That is not accidental. It is the architecture.

% ===========================================================================
\section{Measurement Backbone --- Completed Components}
% ===========================================================================

\begin{table}[h!]
\centering
\begin{tabularx}{\linewidth}{lX}
\toprule
\textbf{Component} & \textbf{Role} \\
\midrule
\texttt{vsepr::Vec3} authoritative state vector & Single canonical position/velocity type \\
Eigen bridge (isolated at analysis boundaries)  & Matrix ops without contaminating core state \\
Kabsch / RMSD analysis                          & Rotation-invariant structural residual \\
Stationarity / settling gates                   & Convergence check before inference \\
Translation / rotation / deformation distinction& Separates motion types from structural change \\
\bottomrule
\end{tabularx}
\caption{Beta-6 measurement infrastructure.}
\end{table}

\textbf{Key validation result:} Kabsch RMSD ignores rigid translation and rotation; it responds
to true non-rigid deformation. Mean displacement captures raw motion. Structural residual
captures reference-site mismatch. This gives the project a defensible way to separate
\emph{moved}, \emph{rotated}, \emph{deformed}, \emph{relaxed}, and \emph{destabilized} ---
a distinction required before any macro inference can be credible.

% ===========================================================================
\section{Task 9 --- Crystal Imperfection Emergence \hfill\statpass}
% ===========================================================================

Crystal imperfection detection validated across SC, FCC, and BCC lattices:

\begin{table}[h!]
\centering
\begin{tabularx}{\linewidth}{lX}
\toprule
\textbf{Variant} & \textbf{Detected behavior} \\
\midrule
ideal            & Low RMSD, no occupancy mismatch, no excess atoms \\
interstitial     & $N_\text{excess} = +1$ \\
vacancy          & $N_\text{excess} = -1$, occupancy mismatch \\
substitutional   & \texttt{structure\_class = ideal}, \texttt{identity\_class = substitutional\_site} \\
thermal\_jitter  & Displacement disorder, no false count defect \\
localized\_impulse & Local relaxation / disruption \\
multi\_site      & Broad disorder / instability \\
\bottomrule
\end{tabularx}
\caption{Crystal imperfection variants and detected behavior.}
\end{table}

\textbf{Primary achievement:} Substitutional defects are separated into structural and identity
channels. The analyzer correctly reports:

\begin{lstlisting}
structure_class = ideal_crystal
identity_class  = substitutional_site
\end{lstlisting}

A substitution can be geometrically clean while identity-defective. The code does not panic and
call everything unstable.

\textbf{Recommended future refinement (beta-7):} Expand classifier vocabulary to include
\texttt{thermal\_disorder}, \texttt{localized\_relaxation}, \texttt{broad\_disorder},
\texttt{atom\_excess}, \texttt{atom\_deficit}, \texttt{substitutional\_site}. The numerical
detectors are already stronger than the current labels.

% ===========================================================================
\section{Task 10 --- Surface Interaction Analysis \hfill\statpass}
% ===========================================================================

Surface interaction behavior validated from trajectory only across six cases:

\begin{table}[h!]
\centering
\begin{tabularx}{\linewidth}{lX}
\toprule
\textbf{Case} & \textbf{Observed classification} \\
\midrule
low normal velocity    & \texttt{no\_interaction} \\
medium normal velocity & \texttt{no\_interaction} \\
grazing                & \texttt{no\_interaction} \\
angled impact          & \texttt{no\_interaction} \\
high normal velocity   & \texttt{temporary\_residence} / \texttt{contact} \\
stream injection       & \texttt{embedded\_event} \\
\bottomrule
\end{tabularx}
\caption{Surface interaction classification results.}
\end{table}

\textbf{Key result:} The stream case detected embedded behavior without storing
\texttt{surface\_damage} or \texttt{embedded} as state truth. Classification is derived entirely
from trajectory displacement and contact duration.

\textbf{Recommended future refinement (beta-7):} Classifier priority should enforce event
hierarchy:
\begin{lstlisting}
embedded_event > ejection_event > temporary_residence > contact_no_capture > no_interaction
\end{lstlisting}

% ===========================================================================
\section{Task 12 / 12B --- Diffusion and Macro Transport Inference}
% ===========================================================================

\textbf{New file:} \texttt{src/analysis/diffusion\_analysis.hpp}

\subsection{Architecture}

\begin{table}[h!]
\centering
\begin{tabularx}{\linewidth}{lX}
\toprule
\textbf{Type} & \textbf{Purpose} \\
\midrule
\texttt{DiffusionRecord}        & Per-frame: MSD total/per-axis, path length, tortuosity, jump count, residence time, neighbor exchange, KE proxy, energy drift, motion region \\
\texttt{DiffusionTracker}       & Feeds position/velocity frames; computes all trajectory-derived metrics \\
\texttt{TransportInference}     & Per-case macro transport inference table (analysis-only) \\
\texttt{TransportAnalyzer}      & Fits MSD slope via linear least-squares; classifies transport regime \\
\texttt{annotate\_defect\_ratios}   & Computes $D_\text{ratio\_vs\_ideal}$ by comparing defect cases to ideal reference \\
\texttt{ActivationTrend::fit}   & Fits $\ln(D)$ vs $1/KE_\text{proxy}$ over multi-temperature cases \\
\texttt{annotate\_formation\_context} & Stamps formation-energy proxy and context label (analysis-only) \\
\bottomrule
\end{tabularx}
\caption{Diffusion analysis module architecture.}
\end{table}

\subsection{Measured Quantities (per frame)}

MSD total; MSD by axis ($x$, $y$, $z$); net displacement; path length; tortuosity proxy; jump
count (site-to-site transitions); residence time; neighbor exchange count; kinetic energy proxy;
energy drift (relative to baseline); motion region (surface / bulk).

\subsection{Inferred Macro Transport Quantities (per case, analysis-only)}

$D_\text{eff}$ (\AA$^2$/fs from MSD slope / $2d$); $D_x$, $D_y$, $D_z$; anisotropy ratio
$= \max(D) / \min_\text{nonzero}(D)$; jump rate; mean residence time; transport class;
anisotropy class; $D_\text{fit quality}$ ($R^2$); $D_\text{ratio\_vs\_ideal}$;
$E_\text{ref}$, $E_\text{atom}$, $\Delta U$ proxy; formation context.

\subsection{Test Results --- 12 Diffusion Cases}

All cases: NVE Lennard-Jones Argon ($Z=18$, $\sigma=3.4$\,\AA, $\varepsilon=0.238$\,kcal/mol,
$m=39.948$\,amu, $\Delta t = 2\times10^{-3}$\,fs, $r_c=10$\,\AA).

\subsubsection{SC 3$\times$3$\times$3 Crystal Cases (12\,000 steps, seed = 77)}

\begin{table}[h!]
\centering
\small
\begin{tabular}{lccccc}
\toprule
\textbf{case\_id} & \textbf{dim} & \textbf{MSD slope (\AA$^2$/fs)} & \textbf{$D_\text{eff}$ (\AA$^2$/fs)} & \textbf{$D_\text{ratio}$} & \textbf{$\Delta U$ (kcal/mol)} \\
\midrule
case\_ideal        & 3 & 0.039578 & 0.006596 & 1.000  & 0.0000  \\
case\_vacancy      & 3 & 0.239472 & 0.039912 & 6.051  & +0.0603 \\
case\_interstitial & 3 & 0.252350 & 0.042058 & 6.376  & +34.631 \\
case\_mixed\_defect& 3 & 0.592508 & 0.098751 & 14.971 & +0.1057 \\
\bottomrule
\end{tabular}
\caption{SC defect diffusion cases.}
\end{table}

\subsubsection{Alternative Crystal Structures}

\begin{table}[h!]
\centering
\small
\begin{tabular}{lccccc}
\toprule
\textbf{case\_id} & \textbf{dim} & \textbf{$D_\text{eff}$ (\AA$^2$/fs)} & \textbf{$E$/atom (kcal/mol)} & \textbf{$\Delta U$} & \textbf{context} \\
\midrule
case\_fcc ($a=5.26$\,\AA) & 3 & 0.030353 & $-1.2530$ & $-0.5883$ & ideal\_reference \\
case\_bcc ($a=3.30$\,\AA) & 3 & 0.025502 & $+11.2433$ & $+11.9080$ & ideal\_reference \\
\bottomrule
\end{tabular}
\caption{FCC and BCC structural comparison cases.}
\end{table}

FCC is lower energy per atom than SC reference ($-0.59$\,kcal/mol), consistent with
LJ-parameterized Ar preferring close-packed structures. BCC at $a=3.30$\,\AA\ is repulsive
($+11.9$\,kcal/mol), correctly flagging over-compressed packing.

\subsubsection{4-Point Thermal Activation Sweep (20\,000 steps each)}

\begin{table}[h!]
\centering
\begin{tabular}{lccc}
\toprule
\textbf{case\_id} & \textbf{$\sigma_\text{jitter}$} & \textbf{$D_\text{eff}$ (\AA$^2$/fs)} & \textbf{$D_\text{ratio\_vs\_ideal}$} \\
\midrule
case\_thermal\_lo  & 0.05 & 0.061358  & $9.30\times$    \\
case\_thermal\_med & 0.12 & 0.367806  & $55.76\times$   \\
case\_thermal\_hi  & 0.30 & 2.385689  & $361.67\times$  \\
case\_thermal\_vhi & 0.55 & 8.184049  & $1240.70\times$ \\
\bottomrule
\end{tabular}
\caption{Thermal sweep: $D_\text{eff}$ spans 130$\times$ across four cases.}
\end{table}

$D_\text{eff}$ spans $130\times$ across the thermal sweep. The activation trend classifier
returns \texttt{insufficient\_data} because all NVE cases flag \texttt{invalid\_energy\_drift}
--- no thermostat, no PBC, energy drifts. This is correct. The flag gates the Arrhenius fit
appropriately.

\subsubsection{Reduced-Dimension Diffusion}

\begin{table}[h!]
\centering
\begin{tabular}{lccl}
\toprule
\textbf{case\_id} & \textbf{dim} & \textbf{$D_\text{eff}$ (\AA$^2$/fs)} & \textbf{anisotropy\_class} \\
\midrule
case\_surface (SC $5\times5\times1$ slab) & 2 & 0.231231 & \texttt{anisotropic\_transport} \\
case\_channel\_1d (SC $1\times1\times20$ rod) & 1 & 1.101889 & \texttt{weakly\_anisotropic} \\
\bottomrule
\end{tabular}
\caption{Reduced-dimension diffusion cases.}
\end{table}

Surface case: $D_z = 0.233$ vs $D_x = 0.112$ / $D_y = 0.117$. The out-of-plane axis shows
higher apparent diffusivity, consistent with the unconstrained $z$-direction in a flat slab
with no PBC.

\subsection{Per-Axis Anisotropy Summary}

\begin{table}[h!]
\centering
\small
\begin{tabular}{lccccc}
\toprule
\textbf{case\_id} & \textbf{$D_x$} & \textbf{$D_y$} & \textbf{$D_z$} & \textbf{aniso ratio} & \textbf{class} \\
\midrule
case\_ideal         & 0.00418 & 0.00813 & 0.00749 & 1.946 & weakly\_anisotropic \\
case\_vacancy       & 0.02494 & 0.04679 & 0.04800 & 1.925 & weakly\_anisotropic \\
case\_interstitial  & 0.03366 & 0.04761 & 0.04491 & 1.414 & weakly\_anisotropic \\
case\_mixed\_defect & 0.06208 & 0.11249 & 0.12169 & 1.960 & weakly\_anisotropic \\
case\_fcc           & 0.02812 & 0.02735 & 0.03559 & 1.302 & weakly\_anisotropic \\
case\_bcc           & 0.03115 & 0.02751 & 0.01786 & 1.744 & weakly\_anisotropic \\
case\_surface       & 0.11209 & 0.11688 & 0.23349 & 2.083 & \textbf{anisotropic\_transport} \\
case\_channel\_1d   & 0.27187 & 0.45973 & 0.37029 & 1.691 & weakly\_anisotropic \\
\bottomrule
\end{tabular}
\caption{Per-axis diffusivity and anisotropy classification. All values in \AA$^2$/fs.}
\end{table}

% ===========================================================================
\section{Task 13 / 13B --- Packing and Macro Property Inference}
% ===========================================================================

\textbf{New file:} \texttt{src/analysis/packing\_analysis.hpp}

\subsection{Architecture}

\begin{table}[h!]
\centering
\begin{tabularx}{\linewidth}{lX}
\toprule
\textbf{Type} & \textbf{Purpose} \\
\midrule
\texttt{PackingRecord}          & Per-frame: packing fraction, coordination, largest cluster, wall proxy, contact persistence, void connectivity, settling flag, energy drift \\
\texttt{PackingTracker}         & $O(N^2)$ contact network from bead geometry; BFS for largest connected cluster \\
\texttt{PackingInference}       & Per-case macro packing table (analysis-only) \\
\texttt{PackingAnalyzer}        & Infers all macro properties from geometry/contact evolution \\
\texttt{compute\_compressibility} & Static; takes initial + compressed snapshots, derives $\Delta U / \Delta V$ proxies \\
\bottomrule
\end{tabularx}
\caption{Packing analysis module architecture.}
\end{table}

\subsection{Inferred Macro Packing Quantities (per case, analysis-only)}

Bulk density proxy (occupied volume / container volume); bulk density inferred (mass-normalized);
porosity $= 1 - \phi$; mean coordination; largest cluster fraction; load-bearing network score;
permeability proxy (Kozeny--Carman-style void/tortuosity estimate); sintering readiness proxy;
packing macro class; compressibility proxy; densification rate proxy; contact growth rate.

\subsection{Test Results --- 4 Packing Cases}

Container: $30 \times 30 \times 30$\,\AA$^3$ (27\,000\,\AA$^3$), bead radius $R = 1.7$\,\AA.

\begin{table}[h!]
\centering
\small
\begin{tabular}{lcccccccl}
\toprule
\textbf{case\_id} & \textbf{N} & \textbf{$\phi$} & \textbf{porosity} & \textbf{$\bar{z}$} & \textbf{cluster} & \textbf{perm.} & \textbf{sinter.} & \textbf{class} \\
\midrule
pack\_loose      &  50 & 0.0381 & 0.9619 & 0.56  & 0.060 & 612.76 & 0.055 & loose\_powder \\
pack\_settled    & 343 & 0.2614 & 0.7386 & 0.018 & 0.006 & 5.894  & 0.001 & loose\_powder \\
pack\_compressed & 343 & 0.2614 & 0.7386 & 5.125 & 1.000 & 5.894  & 0.933 & sintering\_candidate \\
pack\_jammed     & 512 & 0.3902 & 0.6098 & 5.211 & 1.000 & 1.489  & 0.933 & sintering\_candidate \\
\bottomrule
\end{tabular}
\caption{Packing inference results. $\phi$ = packing fraction; $\bar{z}$ = mean coordination; cluster = largest cluster fraction; perm. = permeability proxy; sinter. = sintering readiness proxy.}
\end{table}

\textbf{Compressibility note:} \texttt{densification\_rate = 0.000} (geometry unchanged
between initial/compressed frames in this test); \texttt{contact\_growth\_rate = 5.143}
contacts/frame. The compressibility proxy requires a physical pressure model to become
non-trivial --- labeled as a proxy because scientific honesty is cheaper than retractions.

% ===========================================================================
\section{Group 24 Test Binary}
% ===========================================================================

\begin{tabular}{ll}
\textbf{File}        & \texttt{tests/test\_macro\_inference.cpp} \\
\textbf{CMake group} & 24 \\
\textbf{Target}      & \texttt{test\_macro\_inference} \\
\end{tabular}

\vspace{0.8em}
\textbf{Coverage:}
\begin{itemize}[noitemsep]
	\item 12 diffusion cases (SC ideal, vacancy, interstitial, mixed-defect, FCC, BCC, thermal $\times4$, surface 2D, channel 1D)
	\item 4 packing cases (loose, settled, compressed, jammed)
	\item Formation-context annotation per diffusion case
	\item Per-axis anisotropy audit table
	\item Formation coupling table ($E_\text{ref}$/atom, $E$/atom, $\Delta U$, $D_\text{eff}$)
	\item xyzFull audit assertions on every run
\end{itemize}

\textbf{All pass conditions satisfied:}
\begin{itemize}[noitemsep]
	\item $D_\text{eff}$ computed from MSD slope only \checkmark
	\item Porosity computed from measured packing fraction only \checkmark
	\item No inferred property written to State \checkmark
	\item Energy drift cases correctly flagged \checkmark
	\item Dimensional reduction (2D, 1D) correctly applied \checkmark
	\item Formation context labels correct for all 12 cases \checkmark
	\item Thermal ordering: $D_\text{vhi} \ge D_\text{lo}$ ($130\times$ range confirmed) \checkmark
\end{itemize}

% ===========================================================================
\section{Formation Codebase --- Existing Infrastructure Survey}
% ===========================================================================

The following formation infrastructure exists and is ready for beta-7 wiring:

\begin{table}[h!]
\centering
\begin{tabularx}{\linewidth}{llX}
\toprule
\textbf{Layer} & \textbf{File} & \textbf{Key symbols} \\
\midrule
Formation loop (MD) & \texttt{src/cli/actions\_form.cpp} & \texttt{LangevinDynamics::integrate()}, T-schedule, \texttt{write\_snapshot}, \texttt{formation.log} \\
Formation priors & \texttt{apps/phase4\_formation\_priors.cpp} & crystal preset $\to$ \texttt{UnitCell} $\to$ \texttt{to\_state()} $\to$ single-point energy \\
Thermodynamics & \texttt{include/thermo/thermodynamics.hpp} & \texttt{ThermoData \{H\_f, S, G\_f, Cp\}}, \texttt{ThermoDatabase}, \texttt{GibbsCalculator::calculate()}, \texttt{estimate\_H\_formation()} \\
Fingerprinting & \texttt{atomistic/classify/fingerprints.hpp} & \texttt{ProtoFingerprint}, \texttt{DefectFingerprint}, \texttt{weisfeiler\_lehman\_hash()} \\
Structure clustering & \texttt{atomistic/classify/cluster.hpp} & \texttt{cluster\_by\_proto()}, \texttt{cluster\_by\_defect()}, polymorph/isomorph/defect classify \\
Potential energy & \texttt{src/pot/energy.hpp} & \texttt{EnergyResult \{UvdW, UCoul, total\}}, bond/angle/torsion params \\
Formation regime records & \texttt{tests/test\_runners.hpp} & \texttt{FormationRegimeRecord}, \texttt{FormationHistory}, \texttt{build\_regime\_record()}, \texttt{relaxation\_time()} \\
\bottomrule
\end{tabularx}
\caption{Formation infrastructure ready for beta-7 wiring.}
\end{table}

\textbf{Integration path for 12B $\leftrightarrow$ Formation (beta-7 target):}

\begin{lstlisting}
trajectory
  -> DiffusionTracker -> TransportInference
  +  annotate_formation_context(ref_energy_per_atom from UnitCell preset)
  +  ThermoDatabase::get(formula) -> H_f for true delta-H comparison
  +  cluster_by_defect(snapshots) -> microstate classification
  -> formation_energy_proxy bridges D_eff <-> H_f
\end{lstlisting}

The \texttt{formation\_energy\_proxy} field now present in \texttt{TransportInference} is the
first physical bridge between the trajectory analysis layer and the thermodynamic formation
layer. It is a $\Delta U$ proxy, not a true DFT formation enthalpy. It is labeled accordingly.

% ===========================================================================
\section{Known Limitations and Honest Notes}
% ===========================================================================

\subsection{Energy Drift in NVE Lennard-Jones Tests}

All 12 diffusion cases flag \texttt{invalid\_energy\_drift} and
\texttt{valid\_for\_macro\_inference = no}. This is correct. NVE integration without PBC or
thermostat on a small SC crystal drifts. The flag works. $D_\text{eff}$ values are still
computed and reported --- they reflect kinetic behavior, not thermodynamic equilibrium.

For physically valid $D_\text{eff}$ values, the system requires a thermostat (Langevin already
exists in \texttt{actions\_form.cpp}) and PBC (supported in \texttt{atomistic::Box}). That
wiring is a beta-7 task.

\subsection{BCC Argon at $a = 3.30$\,\AA}

$\Delta U = +11.9$\,kcal/mol vs SC reference because $a = 3.30$\,\AA\ is well inside the LJ
repulsive wall for Ar ($\sigma = 3.4$\,\AA). The BCC test correctly shows
overcompressed-repulsive behavior. It is not a physical BCC Ar ground state. The label
\texttt{ideal\_reference} refers to the geometric reference classification, not a claim of
thermodynamic stability.

\subsection{Interstitial $\Delta U = +34.6$\,kcal/mol}

The interstitial atom is placed at an off-lattice position in the repulsive core of its
neighbors. The large $\Delta U$ correctly reflects that a naive geometric interstitial injection
is high-energy. A FIRE-relaxed interstitial would show a smaller and physically meaningful
formation energy. FIRE relaxation is available in the existing integrator stack and is a beta-7
enhancement target.

\subsection{Activation Trend Returns \texttt{insufficient\_data}}

The 4-point thermal sweep produces a $130\times$ $D$ range but all cases are
\texttt{invalid\_energy\_drift}. The Arrhenius fit gate correctly excludes them. To get a real
activation trend, the cases must be thermostatted. The infrastructure is correct. The physics
input is not yet calibrated.

% ===========================================================================
\section{Beta Release Definitions}
% ===========================================================================

\begin{table}[h!]
\centering
\begin{tabularx}{\linewidth}{llX}
\toprule
\textbf{Version} & \textbf{Label} & \textbf{What ships} \\
\midrule
v4.1.0 & Stable foundation & Frozen known-good baseline \\
\textbf{v5.0.0-beta.6} & \textbf{Formation and Premacro Analysis Foundation} & Formation loop, thermodynamics, fingerprinting, clustering, crystal imperfection, surface analysis, diffusion inference, packing inference, formation-context annotation \\
v5.0.0-beta.7 & Formation/Analysis Unification and Report Freeze & Connect formation outputs $\to$ fingerprint $\to$ cluster $\to$ analyze $\to$ report \\
v5.0.0-beta.8 & Coupled Premacro-to-Macro Workflow Prototype & Material Candidate Card, imperfection $\to$ diffusion workflow, packing $\to$ macro workflow, candidate-level screening \\
\bottomrule
\end{tabularx}
\caption{Beta release roadmap.}
\end{table}

% ===========================================================================
\section{Executive Summary}
% ===========================================================================

Beta-6 completes the instrumentation layer of VSEPR-SIM v5. The project now has:

\begin{itemize}
	\item a formation engine skeleton (formation loop, priors, thermodynamics, fingerprinting, clustering, regime records)
	\item a premacro measurement backbone (Vec3 authority, Eigen bridge, Kabsch/RMSD, stationarity, crystal imperfection, surface interaction)
	\item trajectory-to-inference pipelines for diffusion and packing
	\item analysis-only macro property inference with explicit state-truth separation
	\item formation-context annotation bridging $\Delta U$ proxies to transport behavior
	\item 12 validated diffusion cases and 4 validated packing cases with exact numerical outputs
\end{itemize}

\begin{docnote}
\textbf{The architecture rule held.} No inferred macro property is stored in State.\\[0.4em]
Beta-6 is the instrumentation release. Beta-7 wires the instruments together. Beta-8 uses
them for screening.
\end{docnote}

The formation codebase survey printed in \texttt{test\_macro\_inference} output documents the
precise integration path forward.

\vfill
\noindent\rule{\linewidth}{0.4pt}\\[0.3em]
{\small\color{gray}
Report generated from live test run. All numerical values confirmed reproducible across multiple
executions.\\
Branch: \texttt{4.0-legacy-beta} \quad|\quad Commit: v5.0.0-beta.6 --- Formation and Premacro
Analysis Foundation
}

\end{document}
